\begin{formulaBox}{Einfache lineare Regression}{ML I}
\begin{multicols}{2}
\subsection*{Regressionsgleichung}
\begin{align*}
\hat{y}_i &= b_0+b_1\cdot x_{i}\\
b_1 &= \frac{s_{x, y}}{s_x^2}
     = r_{x, y} \frac{s_y}{s_x}\\
b_0 &= \bar{y}-b_1 \bar{x}\\
y_i &= \hat{y}_i + e_i\\
\beta_1 &= b_1\frac{s_x}{s_y}=r_{x, y}
\end{align*}

\subsection*{Messwertzerlegung}
$$y_i-\bar{y}=(y_i-\hat{y}_i)+(\hat{y}_i-\bar{y})$$

\subsection*{Quadratsummen}
\begin{align*}
\mathrm{QS}_y &= \sum\limits_{i=1}^{n}(y_i-\bar{y})^2\\
\mathrm{QS}_e &= \sum\limits_{i=1}^{n}(y_i-\hat{y}_i)^2\\
\mathrm{QS}_{\hat{y}} &= \sum\limits_{i=1}^{n}(\hat{y}_i-\bar{y})^2\\
\mathrm{QS}_y &= \mathrm{QS}_{\hat{y}}+\mathrm{QS}_e
\end{align*}

\subsection*{Varianzzerlegung}
\begin{align*}
s_y^2 &= \frac{\mathrm{QS}_y}{n}, &
s_e^2 &= \frac{\mathrm{QS}_e}{n}, &
s_{\hat{y}}^2 &= \frac{\mathrm{QS}_{\hat{y}}}{n}\\
s_y^2 &= s_{\hat{y}}^2+s_e^2
\end{align*}

\subsection*{Determinationskoeffizient} 
$$R^2=\frac{s_{\hat{y}}^2}{s_y^2}=\frac{\mathrm{QS}_{\hat{y}}}{\mathrm{QS}_y}=1-\frac{\mathrm{QS}_e}{\mathrm{QS}_y}=r_{x, y}^2=\beta_1^2$$
$$R^2_{\mathrm{adj}}=1-(1-R^2)\cdot\frac{n-1}{n-2} $$

\subsection*{Standardschätzfehler}
$$s_e= \sqrt{\frac{\mathrm{QS}_e}{n}} = s_y\cdot\sqrt{1-R^2}$$
$$\hat{\sigma}_e = \sqrt{\frac{\mathrm{QS}_e}{n-2}}
= \hat{\sigma}_y\cdot\sqrt{1-R^2_{\mathrm{adj}}}$$

\end{multicols}
\end{formulaBox}

\begin{formulaBox}{Multiple Regression}{ML II}

 \begin{multicols}{2}
In den folgenden Formeln wird zur Veranschaulichung der Fall mit zwei Prädiktoren dargestellt. 
Für mehr als zwei Prädiktoren lassen sich die Regressionsgleichung, die Beziehung zwischen $b$ und $\beta$, 
die Berechnung von $b_0$ sowie die Formeln für den Standardschätzfehler analog erweitern. 
Die Berechnung der $\beta$-Gewichte und von $R^2$ erfordert jedoch eine Matrixschreibweise, die nicht überall unterichtet wird.

    \subsection*{Regressionsgleichung}
für Originalwerte: $$\hat{y}_i=b_0+b_1 \cdot x_{1i} + b_2 \cdot x_{2i}$$
wobei: 
\begin{align*}
b_1&=\beta_1 \cdot \frac{s_y}{s_{x_{1}}}\\
b_2&=\beta_2 \cdot \frac{s_y}{s_{x_{2}}}\\
b_0&=\bar{y}-b_1 \cdot \bar{x}_{1}-b_2 \cdot \bar{x}_2\\
\end{align*}

für z-standardisierte Werte: $$\hat{z}_{y_i}=\beta_1\cdot z_{x_{1i}} + \beta_2\cdot z_{x_{2i}}$$ \\
wobei: $$\beta_1=\frac{r_{y,x_1}-r_{y,x_2}\cdot r_{x_1,x_2}}{1-r^2_{x_1,x_2}}$$ \ $$\beta_2=\frac{r_{y,x_2}-r_{y,x_1} \cdot r_{x_1,x_2}}{1-r^2_{x_1,x_2}}$$\\

\columnbreak
    \subsection*{Gütemaße}

Multipler Determinationskoeffizient
$$R^2=\beta_1 \cdot r_{y,x_1} + \beta_2 \cdot r_{y,x_2}$$
$$R^2_{\mathrm{adj}}=1-(1-R^2)\cdot\frac{n-1}{n-k-1} $$
Multipler Korrelationskoeffizient
$$R=\sqrt{R^2}$$
Standardschätzfehler für Originalwerte

$$s_e=\sqrt{\frac{\sum\limits_{i=1}^n(y_i-\hat{y}_i)^2}{n}} = s_y \sqrt{1-R^2}$$

$$\hat{\sigma}_e = \sqrt{\frac{\sum\limits_{i=1}^n(y_i-\hat{y}_i)^2}{n-k-1}} = \hat{\sigma}_y \cdot \sqrt{1-R^2_{\mathrm{adj}}}$$

Standardschätzfehler für $z$-standardisierte Werte
$$\hat{\sigma}_e=\sqrt{1-R^2_{\mathrm{adj}}}$$

\end{multicols}
\end{formulaBox}
